\newpage
\phantomsection
\label{prf:between-any-two-rationals-is-a-rational}
\LRAProofFor{thm:between-any-two-rationals-is-a-rational}

\begin{remark*}[Return]
\hyperref[thm:between-any-two-rationals-is-a-rational]{Return to Theorem}
\end{remark*}

\begin{theorem*}[Between Any Two Rational Numbers Lies Another Rational]
If $r,s\in\mathbb{Q}$ and $r<s$, then there exists $t\in\mathbb{Q}$ such that
\[
r<t<s.
\]
\end{theorem*}

\begin{proof}[Professional Standard Proof]
\LRAProofBodyStart
Let $r,s\in\mathbb{Q}$ with
\[
r<s.
\]
Set
\[
t:=\frac{r+s}{2}.
\]
Since $\mathbb{Q}$ is closed under addition and division by the nonzero
rational number $2$, we have
\[
t\in\mathbb{Q}.
\]

It remains to show that $t$ lies strictly between $r$ and $s$.

Because $r<s$, adding $r$ to both sides gives
\[
r+r<r+s.
\]
Dividing by the positive rational number $2$ yields
\[
r<\frac{r+s}{2}=t.
\]

Similarly, from $r<s$, adding $s$ to both sides gives
\[
r+s<s+s.
\]
Dividing by the positive rational number $2$ yields
\[
t=\frac{r+s}{2}<s.
\]

Therefore
\[
r<t<s,
\]
with $t\in\mathbb{Q}$.  This is exactly the desired conclusion.
\end{proof}

\begin{remark*}[Proof structure]
The proof uses the midpoint construction.  Given two rationals $r<s$, their
average $(r+s)/2$ is again rational by closure of $\mathbb{Q}$ under the
field operations.  The inequalities $r<(r+s)/2$ and $(r+s)/2<s$ come from the
original inequality $r<s$ by adding the same endpoint to both sides and then
dividing by the positive rational number $2$.  So the midpoint provides the
required interior rational.
\end{remark*}

\begin{dependencies}
\begin{itemize}
  \item Basic arithmetic closure of $\mathbb{Q}$ under addition and division by $2$
  \item \hyperref[thm:one-positive]{Theorem~\ref*{thm:one-positive}}
        \textnormal{(to know that $2=1+1$ is positive)}
\end{itemize}
\end{dependencies}

\begin{proof}[Detailed Learning Proof]
\LRAProofBodyStart
We expand the professional proof by following the same sequence of justified moves.

Let $r,s\in\mathbb{Q}$ with
\[
r<s.
\]
Set
\[
t:=\frac{r+s}{2}.
\]
Since $\mathbb{Q}$ is closed under addition and division by the nonzero
rational number $2$, we have
\[
t\in\mathbb{Q}.
\]

It remains to show that $t$ lies strictly between $r$ and $s$.

Because $r<s$, adding $r$ to both sides gives
\[
r+r<r+s.
\]
Dividing by the positive rational number $2$ yields
\[
r<\frac{r+s}{2}=t.
\]

Similarly, from $r<s$, adding $s$ to both sides gives
\[
r+s<s+s.
\]
Dividing by the positive rational number $2$ yields
\[
t=\frac{r+s}{2}<s.
\]

Therefore
\[
r<t<s,
\]
with $t\in\mathbb{Q}$.  This is exactly the desired conclusion.
\end{proof}
